\section{The effective action along the clock}
\label{sec:cjt}

\subsection{The method, and why it carries over}

The results on the paired state, on matter, on temperature and on the
vacuum follow from the effective action for composite operators of Cornwall, Jackiw and Tomboulis (CJT)
\cite{CJT1974}: a functional of the expectation value $\phi$ of the Vorton-Scalar and of the full
propagators, whose two-particle-irreducible part is truncated at two loops. Its stationarity
conditions are the equation for the fermion mass in matter and the gap equation. On the deflating
spacetime, and for fields that do not depend on the extra coordinates, nothing in this construction
changes, for two reasons established above. The weighted fermion field obeys the flat equations, so
that its propagator is the flat one. And the interaction between the fermions is, at each moment,
the exchange of the two modes of the scalar with the coupling of that moment and with the unshifted
range (Sec.~\ref{sec:scalarmodes}).

The one new term, the shift $-\gsum^{2}/4$ of the squares of the boson masses in
Eq.~(\ref{eq:couplings}), is smaller than every other scale of the calculation by the square of the
ratio of the rate of deflation to a particle mass. For the largest rate that observation allows
(Sec.~\ref{sec:cosmo}) and a mass of 1~keV that ratio is below $10^{-36}$. The term is neglected in
this section; this is the only approximation made here beyond the truncation itself. The test
of the truncation on a solvable model in zero dimensions, which does not involve the spacetime, was
executed again: the two-loop truncation has a solution for every coupling from 0.01 to 30 and stays
within 10.7 per cent of the exact propagator.

Every result of this section is a statement about a \emph{sequence of equilibrium states}. Whether a
system can follow the sequence was examined, for the paired state, in Sec.~\ref{sec:quench}.

\subsection{The paired state}

The CJT potential $V[\Delta]$ of the paired state is the effective action per unit volume, measured
from the unpaired state, as a functional of the gap function $\Delta(k)$. The program evaluates it at
13 moments, for two choices of the mass $m_s$ at the Fermi momentum $\kF=0.5\,m$, in the vector
channel (Table~\ref{tab:cjt}, Fig.~\ref{fig:cjt}).

\begin{table*}
\caption{\label{tab:cjt}The paired state along the clock from the CJT potential, for $\kF=0.5\,m$.
$\lambda$: pairing strength. $V$: the potential at the solution, the condensation energy. Ratio: $V$
divided by its weak-coupling value $-\nu N(0)\Delta_F^{2}/4$, with $\nu=8$ internal states and the
density of states $N(0)$. Curvature: the smallest eigenvalue of the matrix of second derivatives of
the potential. Size: the coherence length of the pair in units of $1/\kF$.}
\begin{ruledtabular}
\begin{tabular}{cccccccccc}
$m_s/m$, $\kF/m$ & $\td$ & $\gVS$ & $\lambda$ & $\Delta_F/m$ & $V/m^{4}$ & ratio & curvature & $T_c/\Delta_F$ & size $\times\kF$ \\
\hline
{{ROWS:tbl_cjt}}
\end{tabular}
\end{ruledtabular}
\end{table*}

\emph{The solution of the gap equation is a minimum at every moment.} Newton's method applied to the
potential converges to the gap function that the iteration of the gap equation finds, to
$7.7\times10^{-12}$, and the smallest eigenvalue of the second variation is positive throughout. The
temperature at which the second variation at $\Delta=0$ loses its positivity equals the temperature
at which the linearized gap equation has the eigenvalue 1, to $2.6\times10^{-12}$: two routes to the
critical temperature agree.

\emph{The condensation energy deepens.} It is within 15 per cent of the weak-coupling value wherever
the gap is below $0.1\,\EF$, and in weak coupling its logarithmic slope along the clock is twice that
of the gap, by Eq.~(\ref{eq:growth}):
\begin{equation}
\frac{d\ln|V|}{d\td}=\frac{2}{\lambda}.
\end{equation}

\emph{The pairs shrink.} The coherence length of the first set falls from 25 to 1.0 in units of
$1/\kF$. At its last five moments the gap exceeds $0.1\,\EF$, and at the last the pair is as small
as the distance between the particles. \emph{The deflation carries the paired state from overlapping
pairs toward bound pairs.} On flat spacetime this crossover \cite{NozieresSchmittRink1985,Leggett1975}
is reached by choosing a larger coupling; here it is reached by waiting.

\emph{The universal ratios hold while the coupling is weak.} While the critical temperature rises by
a factor of sixteen, from $6.3\times10^{-3}$ to $0.10$ of the Fermi energy, the ratio of the gap at
zero temperature to the critical temperature stays within 0.019 of the weak-coupling value 1.7639,
and the relative jump of the heat capacity at $T_c$ within 0.048 of 1.4261. The coefficients of the
Ginzburg-Landau energy agree with their weak-coupling formulas to 1.0 per cent while $T_c$ is below
$0.05\,\EF$ and differ by up to 7.9 per cent at $T_c=0.10\,\EF$. The coherence length of the
Ginzburg-Landau energy, $0.1333\,v_F/T_c$ in weak coupling with the Fermi velocity $v_F$, shrinks
from 80 to 5.1 in units of $1/m$: the cores of the vortices of the superfluid shrink as the
deflation proceeds.

\begin{figure*}
\includegraphics[width=\textwidth]{fig_cjt_history}
\caption{\label{fig:cjt}The paired state along the clock, from the CJT potential, for the two sets of
Table~\ref{tab:cjt}. (a) The condensation energy against the coupling of the moment. (b) The
coherence length of the pair in units of $1/\kF$. (c) The condensation energy divided by its
weak-coupling value, against the gap in units of the Fermi kinetic energy: the departure from weak
coupling.}
\end{figure*}

\subsection{Matter}

\emph{Matter of fixed density loses its fermion mass.} In matter the expectation value of the
Vorton-Scalar is reduced by the scalar density of the fermions, and with it the fermion mass in the
matter, $m^{*}$. The reduction depends on the density and on the
single number $C$, and above a critical value of $C$ the mass vanishes: the matter is then in
the \emph{abnormal} state of Lee and Wick \cite{LeeWick1974}, a gas of massless fermions in a region
without expectation value of the scalar. On the deflating spacetime $C\propto e^{\td}$ grows at
fixed density. The critical values are $C=151.892$, $12.971$, $2.556$ and $0.617$ for the Fermi
momenta $\kF=0.2$, $0.5$, $1$ and $2$ in units of $m$; they exceed the estimate
$\pi^{2}m^{2}/4\kF^{2}$, which holds for dense matter, by the factors 2.46, 1.31, 1.04 and 1.00. \emph{Every
piece of matter becomes abnormal at a definite moment, the denser the earlier.}

\emph{Self-bound matter.} Matter is self-bound, that is, its energy per particle at some density
lies below the mass of a free fermion, if $C>3\pi^{2}/8=3.7011$ without the filled negative-energy
states and if $C>3\pi^{2}/5=5.9218$ with them. Once the threshold is passed, the lowest energy per
particle falls monotonically along the clock (Table~\ref{tab:saturation}), from $-0.0192\,m$ at
$C=4$ to $-0.405\,m$ at $C=29.6$. At every one of these moments the matter of lowest energy is in the
abnormal state. The exchange energies of the two modes, which the mean field omits, are 0.44 and
0.51 per cent of the mean-field energy at the first moment and 0.90 and 1.22 per cent at the last
(for $m_s/m=0.3$):
the mean field remains a good approximation along the whole history. A gas of the fermions that is
dilute and unbound at early times thus has, at late times, a state of lower energy: drops of dense
abnormal matter, bound by up to 40 per cent of the rest mass at the last moment of the table.
Whether and how fast a gas reaches that state is a question of kinetics that is not answered here.

\begin{table}
\caption{\label{tab:saturation}Self-bound matter along the clock, in the mean-field approximation
without the negative-energy states, for $C=4\,e^{\td}$: the lowest energy per particle minus the
fermion mass, and the Fermi momentum at which it is reached.}
\begin{ruledtabular}
\begin{tabular}{cccc}
$\td$ & $C$ & $(E/N-m)/m$ & $\kF/m$ \\
\hline
{{ROWS:tbl_saturation|cols=0,1,2,3}}
\end{tabular}
\end{ruledtabular}
\end{table}

\subsection{Temperature}

At high temperature the expectation value of the Vorton-Scalar vanishes and the fermions are
massless. The temperature $T_{1/2}$ at which the expectation value has fallen to one half is of the
order of the four-dimensional vacuum value $\vVS=m/\gVS$, which falls as $e^{-\td/2}$ while the
masses stay fixed. In units of the constant fermion mass, $T_{1/2}$ falls monotonically as the
coupling grows, for each of the three ratios $m_s/m=0.3$, $1$ and $3$ examined; for $m_s/m=1$, from
$2.33\,m$ at $\gVS=0.3$ to $0.37\,m$ at $\gVS=4$ [Fig.~\ref{fig:matter}(d)]. \emph{Matter of a fixed
temperature melts at a definite moment of the deflation.} For $m_s/m=0.3$ no melting temperature
can be given from $C=32$ on: there the one-loop potential of the vacuum turns over so close to the
vacuum value that the Hartree potential ceases to exist before the expectation value has fallen, and
the program marks these moments as not defined.

\subsection{The vacuum}

The one-loop potential of the Vorton-Scalar in vacuum contains the negative contribution of the
fermion loop \cite{ColemanWeinberg1973}, which grows with the fourth power of the field and
eventually wins: beyond some field the potential lies below its value in the vacuum. The larger $C$,
the smaller that field: 24.7 times the vacuum value at $C=1$, and 1.77 times at $C=10$
[Fig.~\ref{fig:matter}(c)]. As the deflation proceeds, the barrier between the vacuum of the model
and the region where the one-loop potential is unbounded below moves in toward the vacuum.

The one-loop potential is not reliable at fields far from the vacuum value, and a potential that
falls at large field signals that the effective theory needs completion there, not that the vacuum
decays. What the result establishes is a trend: every statement of the model that rests on the
vacuum of the Vorton-Scalar becomes less secure with time. The order of this event relative to the
binding of the pair was given in Sec.~\ref{sec:pair}.

\emph{A pair condensate in the vacuum.} The ladder equation with a
momentum cutoff shows that the pair can condense without a medium only if $\aVS$ exceeds a critical value
$\alpha_c$ of order one \cite{Miransky1985}; with the cutoff $100\,m$ the critical values lie between
2.13 and 2.93 for the six mass ratios of Table~\ref{tab:events}. On the deflating spacetime $\aVS$
reaches such a value at a definite moment. But at that moment the coupling is strong, the exchange
of one boson is no longer a controlled approximation, and the vacuum of the model is no longer
protected. A pair condensate in the vacuum is not reached by the deflation within the range of
validity of the model.

\begin{figure*}
\includegraphics[width=0.92\textwidth]{fig_matter_history}
\caption{\label{fig:matter}Matter, vacuum and temperature along the clock. (a) The fermion mass in
matter of three fixed densities against $C\propto e^{\td}$: it drops to zero at the critical value.
(b) The lowest energy per particle of self-bound matter, minus $m$. (c) The field at which the
one-loop potential of the vacuum has its maximum, and the field beyond which it lies below the vacuum
value, in units of the vacuum value. (d) The temperature at which the fermions have lost half of
their mass, against the coupling, for three ratios $m_s/m$.}
\end{figure*}

